% ============================================================
%  Section 5 — Results
%  Measured outcomes only; the methods are defined in Sections 3-4.
%  All figures: XJTU-SY, 5 outer folds, 5 seeds, alpha = 0.10 unless
%  stated; mean +/- std over seeds of the fold-averaged metrics.
%  Source: repo results/09_cvplus (paper_tables.json,
%  per_bearing_report.md) and results/c3_offline.
% ============================================================

\section{Results}
\label{sec:results}

% ------------------------------------------------------------
\subsection{Synthetic Near-Failure Windows}
\label{sec:res_gan}

The synthetic near-failure windows were assessed on four criteria
before entering training: separability from real windows, diversity,
location in a two-dimensional embedding, and per-feature marginal
distributions. Table~\ref{tab:gan_quality} reports the first two for
the generators of the five outer folds, trained on the same
configuration as the inner generators.

\begin{table}[htbp]
\centering
\caption{Quality of the synthetic near-failure windows per fold.
Separability: accuracy of a classifier separating real from synthetic
windows (50\% = indistinguishable). Diversity ratio: mean pairwise
distance among synthetic windows divided by that among real windows
(1 = matched diversity).}
\label{tab:gan_quality}
\footnotesize
\renewcommand{\arraystretch}{1.2}
\begin{tabular}{cccccc}
\toprule
& \multicolumn{5}{c}{\textbf{Fold}} \\
\cmidrule(lr){2-6}
\textbf{Metric} & 1 & 2 & 3 & 4 & 5 \\
\midrule
Separability (\%) & 67.7 & 59.1 & 61.5 & 70.2 & 64.7 \\
Diversity ratio   & 0.82 & 0.78 & 0.81 & 0.78 & 0.87 \\
\bottomrule
\end{tabular}
\end{table}

Separability stays between 59\% and 70\%, approaching chance in the
folds richest in near-failure data, and the diversity ratio between
0.78 and 0.87, so the generator reproduces the spread of the real data
without mode collapse. In the embedding of Fig.~\ref{fig:tsne}, the
synthetic windows occupy the same regions as the real ones in every
fold, including the small isolated sub-populations, and form no
separate cluster. The per-feature distributions of
Fig.~\ref{fig:feature_dist} coincide for all twenty features,
including the bimodal shape of the amplitude features near failure.

\begin{figure}[htbp]
  \centering
  \includegraphics[width=\textwidth]{figures/04_tsne_real_vs_synthetic.png}
  \caption{t-SNE embedding of real (blue) and synthetic (red)
  near-failure windows for each fold.}
  \label{fig:tsne}
\end{figure}

\begin{figure}[htbp]
  \centering
  \includegraphics[width=\columnwidth]{figures/04_feature_distributions.png}
  \caption{Per-feature distributions of real and synthetic near-failure
  windows for a representative fold.}
  \label{fig:feature_dist}
\end{figure}

Two labelling schemes for the synthetic windows were compared in an
exploratory sweep (Table~\ref{tab:aug_sweep}): the label of the nearest
real near-failure window in feature space, and the sampling scheme of
Eq.~\eqref{eq:syn_label}. The nearest-neighbour scheme shifted the
synthetic labels upward (mean 0.27 against 0.24 for real windows),
pushing the predictor toward overestimating the remaining life. The
sampling scheme matched the real mean within 0.01 in every fold and
gave a lower error at every ratio, with its minimum at $r = 0.75$. A
control in which real near-failure windows were duplicated instead of
synthesised gave no improvement (near-failure RMSE 0.247 at ratio 0.5
and 0.250 at ratio 1.0), which indicates that the predictor was not
limited by the number of near-failure windows as such.

\begin{table}[htbp]
\centering
\caption{Exploratory sweep (single run per setting) of near-failure
RMSE by labelling scheme and augmentation ratio.}
\label{tab:aug_sweep}
\footnotesize
\renewcommand{\arraystretch}{1.2}
\begin{tabular}{lcccc}
\toprule
\textbf{Scheme} & \textbf{0.25} & \textbf{0.50} & \textbf{0.75} & \textbf{1.00} \\
\midrule
Sampling          & 0.272 & 0.247 & \textbf{0.231} & 0.232 \\
Nearest-neighbour & 0.266 & 0.265 & 0.239 & 0.243 \\
\bottomrule
\end{tabular}
\end{table}

% ------------------------------------------------------------
\subsection{Point Accuracy and Stage Recognition}
\label{sec:res_accuracy}

Table~\ref{tab:rmse} reports the RMSE of the four configurations by
stage. Multi-task learning is the component that improves accuracy.
It lowers the near-failure RMSE from 0.225 to 0.216, and to 0.206 when
combined with augmentation, the lowest and least variable value. The
effect is larger over the whole life: the overall RMSE falls from 0.385
to 0.297 with C2 and to 0.291 with both components. Most of this gain
comes from the healthy windows, whose error is the largest and most
variable in the baseline (0.392 $\pm$ 0.290) and drops to about 0.26
with C2. Augmentation alone leaves the near-failure RMSE unchanged
(0.226) and lowers the overall RMSE from 0.385 to 0.347, a difference
within the seed-to-seed variability of that metric.

\begin{table}[htbp]
\centering
\caption{RMSE (normalised RUL, per bearing) by stage, mean $\pm$
standard deviation over five seeds. Point estimate: mean of the four
inner models.}
\label{tab:rmse}
\footnotesize
\renewcommand{\arraystretch}{1.2}
\setlength{\tabcolsep}{4pt}
\begin{tabular}{lccccc}
\toprule
\textbf{Config.} & \textbf{Healthy} & \textbf{Early} &
\textbf{Near-failure} & \textbf{Post-FPT} & \textbf{All} \\
\midrule
Baseline & $0.392 \pm 0.290$ & $0.360 \pm 0.120$ & $0.225 \pm 0.041$ & $0.339 \pm 0.068$ & $0.385 \pm 0.161$ \\
+GAN     & $0.337 \pm 0.226$ & $0.329 \pm 0.084$ & $0.226 \pm 0.038$ & $0.313 \pm 0.042$ & $0.347 \pm 0.121$ \\
+MT      & $0.261 \pm 0.173$ & $0.280 \pm 0.068$ & $0.216 \pm 0.035$ & $0.284 \pm 0.045$ & $0.297 \pm 0.101$ \\
+GAN+MT  & $0.266 \pm 0.166$ & $0.280 \pm 0.075$ & $\mathbf{0.206 \pm 0.026}$ & $\mathbf{0.275 \pm 0.041}$ & $\mathbf{0.291 \pm 0.098}$ \\
\bottomrule
\end{tabular}
\end{table}

The stage on which C3 is conditioned is recognised with an accuracy of
$0.722 \pm 0.021$ by the stage network of the configurations without
augmentation, and $0.706 \pm 0.018$ by the stage head of the full
framework. Near-failure windows are found reliably, with a recall of
$0.887 \pm 0.018$ and $0.896 \pm 0.034$ respectively, but the
precision is lower ($0.730 \pm 0.025$ and $0.674 \pm 0.042$): the
main error is to assign early-degradation windows to the near-failure
stage.

% ------------------------------------------------------------
\subsection{Calibration on Two Bearings versus Cross-Conformal Calibration}
\label{sec:res_calibration}

To measure the effect of the calibration scheme, the configurations
were also calibrated by split conformal prediction, with one predictor
per fold trained on 9~bearings and two held-out calibration bearings
chosen to contain at least 100 near-failure windows and a stage
composition close to the test set. Table~\ref{tab:calibration}
compares this protocol with CV+ for the baseline and the full
framework.

\begin{table*}[!htbp]
\centering
\caption{Coverage and width at a 90\% target under split conformal
calibration on two bearings and under cross-conformal (CV+)
calibration on twelve bearings. PICP $\pm$ std over seeds / MPIW.}
\label{tab:calibration}
\scriptsize
\renewcommand{\arraystretch}{1.2}
\setlength{\tabcolsep}{4pt}
\begin{tabular}{llcccc}
\toprule
& & \multicolumn{2}{c}{\textbf{Baseline}} & \multicolumn{2}{c}{\textbf{+GAN+MT}} \\
\cmidrule(lr){3-4}\cmidrule(lr){5-6}
\textbf{Calibration} & \textbf{Grouping} & \textbf{All} & \textbf{Near-failure} & \textbf{All} & \textbf{Near-failure} \\
\midrule
\multirow{3}{*}{Split, 2 bearings}
 & Single $q$        & $0.793 \pm 0.069$ / 0.68 & $0.871 \pm 0.105$ / 0.69 & $0.830 \pm 0.050$ / 0.59 & $0.839 \pm 0.080$ / 0.63 \\
 & Predicted stage   & $0.740 \pm 0.022$ / 0.63 & $0.887 \pm 0.046$ / 0.65 & $0.817 \pm 0.047$ / 0.54 & $0.884 \pm 0.046$ / 0.67 \\
 & True stage        & $0.786 \pm 0.052$ / 0.55 & $0.860 \pm 0.027$ / 0.47 & $0.811 \pm 0.034$ / 0.46 & $0.864 \pm 0.025$ / 0.49 \\
\midrule
\multirow{3}{*}{CV+, 12 bearings}
 & Single $q$        & $0.898 \pm 0.087$ / 0.84 & $0.971 \pm 0.029$ / 0.88 & $0.856 \pm 0.043$ / 0.75 & $0.959 \pm 0.024$ / 0.79 \\
 & Predicted stage   & $\mathbf{0.909 \pm 0.040}$ / 0.67 & $0.968 \pm 0.028$ / 0.83 & $\mathbf{0.917 \pm 0.039}$ / 0.60 & $0.954 \pm 0.023$ / 0.84 \\
 & True stage        & $0.898 \pm 0.065$ / 0.76 & $0.905 \pm 0.027$ / 0.54 & $0.880 \pm 0.038$ / 0.69 & $0.910 \pm 0.014$ / 0.55 \\
\bottomrule
\end{tabular}
\end{table*}

Calibrated on two bearings, every version under-covers: overall
coverage lies between 0.74 and 0.83, and even the true stage gives only
0.79--0.81 overall and 0.86 near failure. The intervals are narrow
because they are calibrated on two bearings that do not represent the
test bearings. Under CV+, coverage returns to the nominal level: the
predicted-stage version reaches 0.909 and 0.917 overall, and the
true-stage version 0.905 and 0.910 near failure. The intervals are
wider than under split calibration, which is the expected price of
calibrating against the full variability of the non-test bearings.
All remaining results use CV+.

% ------------------------------------------------------------
\subsection{Interval Quality and Stage Conditioning}
\label{sec:res_intervals}

Table~\ref{tab:intervals} compares the three versions of C3 for the
four configurations.

\begin{table*}[!htbp]
\centering
\caption{CV+ coverage and width at a 90\% target by grouping and
configuration. PICP $\pm$ std over seeds / MPIW (clipped). Stage
subsets are defined by the true stage. The predicted stage comes from
the model's own stage head (+MT, +GAN+MT) or from the separate stage
network (baseline, +GAN).}
\label{tab:intervals}
\scriptsize
\renewcommand{\arraystretch}{1.2}
\setlength{\tabcolsep}{2.5pt}
\begin{tabular}{llccccc}
\toprule
\textbf{Config.} & \textbf{Grouping} & \textbf{Healthy} & \textbf{Early} &
\textbf{Near-failure} & \textbf{Post-FPT} & \textbf{All} \\
\midrule
\multirow{3}{*}{Baseline}
 & Single $q$      & $0.825 \pm 0.153$ / 0.83 & $0.902 \pm 0.080$ / 0.87 & $0.971 \pm 0.029$ / 0.88 & $0.932 \pm 0.057$ / 0.87 & $0.898 \pm 0.087$ / 0.84 \\
 & Predicted stage & $0.900 \pm 0.037$ / 0.62 & $0.819 \pm 0.100$ / 0.77 & $0.968 \pm 0.028$ / 0.83 & $0.885 \pm 0.067$ / 0.80 & $0.909 \pm 0.040$ / 0.67 \\
 & True stage      & $0.827 \pm 0.153$ / 0.80 & $0.944 \pm 0.013$ / 0.88 & $0.905 \pm 0.027$ / 0.54 & $0.925 \pm 0.013$ / 0.73 & $0.898 \pm 0.065$ / 0.76 \\
\midrule
\multirow{3}{*}{+GAN}
 & Single $q$      & $0.846 \pm 0.086$ / 0.86 & $0.901 \pm 0.079$ / 0.88 & $0.972 \pm 0.028$ / 0.90 & $0.932 \pm 0.055$ / 0.89 & $0.905 \pm 0.064$ / 0.86 \\
 & Predicted stage & $0.799 \pm 0.062$ / 0.60 & $0.851 \pm 0.088$ / 0.80 & $0.969 \pm 0.028$ / 0.82 & $0.903 \pm 0.057$ / 0.81 & $0.872 \pm 0.030$ / 0.66 \\
 & True stage      & $0.853 \pm 0.086$ / 0.84 & $0.939 \pm 0.020$ / 0.89 & $0.908 \pm 0.022$ / 0.54 & $0.923 \pm 0.021$ / 0.74 & $0.904 \pm 0.045$ / 0.79 \\
\midrule
\multirow{3}{*}{+MT}
 & Single $q$      & $0.807 \pm 0.100$ / 0.73 & $0.833 \pm 0.100$ / 0.79 & $0.955 \pm 0.044$ / 0.79 & $0.886 \pm 0.075$ / 0.79 & $0.860 \pm 0.076$ / 0.74 \\
 & Predicted stage & $0.931 \pm 0.013$ / 0.53 & $0.759 \pm 0.089$ / 0.67 & $0.960 \pm 0.025$ / 0.80 & $0.847 \pm 0.057$ / 0.73 & $0.909 \pm 0.026$ / 0.60 \\
 & True stage      & $0.807 \pm 0.100$ / 0.69 & $0.911 \pm 0.040$ / 0.86 & $0.893 \pm 0.007$ / 0.58 & $0.898 \pm 0.026$ / 0.73 & $0.871 \pm 0.040$ / 0.69 \\
\midrule
\multirow{3}{*}{+GAN+MT}
 & Single $q$      & $0.793 \pm 0.077$ / 0.74 & $0.823 \pm 0.039$ / 0.79 & $0.959 \pm 0.024$ / 0.79 & $0.882 \pm 0.031$ / 0.79 & $0.856 \pm 0.043$ / 0.75 \\
 & Predicted stage & $0.900 \pm 0.051$ / 0.52 & $0.809 \pm 0.133$ / 0.74 & $0.954 \pm 0.023$ / 0.84 & $0.871 \pm 0.080$ / 0.78 & $\mathbf{0.917 \pm 0.039}$ / \textbf{0.60} \\
 & True stage      & $0.793 \pm 0.077$ / 0.70 & $0.943 \pm 0.022$ / 0.88 & $0.910 \pm 0.014$ / 0.55 & $0.925 \pm 0.020$ / 0.73 & $0.880 \pm 0.038$ / 0.69 \\
\bottomrule
\end{tabular}
\end{table*}

\textit{Effect of stage conditioning.} Conditioning on the predicted
stage gives the narrowest intervals over the whole life in every
configuration, at a coverage equal to or above that of the single
quantile except for +GAN (0.872 against 0.905). For the full framework, the overall width falls from 0.75 to
0.60, about 20\%, while coverage rises from 0.856 to 0.917. The gain
comes from the healthy windows, whose intervals narrow from 0.74 to
0.52 while their coverage rises from 0.79 to 0.90. Within the predicted
stages, on which the guarantee is defined, the full framework covers
$0.941 \pm 0.037$ of the windows predicted healthy,
$0.895 \pm 0.072$ of those predicted in early degradation and
$0.937 \pm 0.035$ of those predicted near failure.\\

\textit{Cost of stage errors.} Measured on the true stages, the
predicted-stage version over-covers the near-failure windows
(0.95--0.97) with intervals of 0.80--0.84, and under-covers the early
degradation windows (0.76--0.85). The true-stage reference shows what
correct stage recognition would give: 0.89--0.91 near failure with
intervals of 0.54--0.58, about 35\% narrower. The difference follows
from the precision of Section~\ref{sec:res_accuracy}: early-degradation
windows predicted as near-failure enter the near-failure calibration
group with larger residuals and widen its quantile, while the
near-failure windows themselves receive more width than they need.\\

\textit{Effect of C1 and C2 on the intervals.} With the predicted
stage, multi-task learning narrows the intervals over the whole life
from 0.67 to 0.60 at unchanged coverage (0.909). Augmentation leaves
the width unchanged (0.66) and lowers the overall coverage of the
single-task model to 0.872, a difference of the order of one seed
standard deviation. Combined with C2, augmentation gives the same
width as C2 alone (0.60) and a coverage of 0.917.\\

\textit{Other confidence levels.} Table~\ref{tab:alpha} reports the
predicted-stage version at four levels. Coverage follows the nominal
level over the whole life for the full framework, from 0.954 at a 95\%
target to 0.882 at 80\%, with the width shrinking accordingly. The
near-failure coverage stays between 0.95 and 0.97 at every level,
consistent with the over-coverage described above, and the post-FPT
coverage remains a few points below target.

\begin{table}[htbp]
\centering
\caption{CV+ with predicted-stage conditioning at four confidence
levels: PICP over all windows, post-FPT and near-failure windows, and
MPIW over all windows.}
\label{tab:alpha}
\footnotesize
\renewcommand{\arraystretch}{1.2}
\setlength{\tabcolsep}{4pt}
\begin{tabular}{lccccc}
\toprule
\textbf{Config.} & \textbf{Target} & \textbf{PICP all} &
\textbf{PICP post-FPT} & \textbf{PICP near-failure} & \textbf{MPIW all} \\
\midrule
\multirow{4}{*}{Baseline}
 & 0.95 & 0.935 & 0.923 & 0.969 & 0.74 \\
 & 0.90 & 0.909 & 0.885 & 0.968 & 0.67 \\
 & 0.85 & 0.890 & 0.854 & 0.966 & 0.62 \\
 & 0.80 & 0.843 & 0.835 & 0.962 & 0.60 \\
\midrule
\multirow{4}{*}{+GAN+MT}
 & 0.95 & 0.954 & 0.908 & 0.959 & 0.66 \\
 & 0.90 & 0.917 & 0.871 & 0.954 & 0.60 \\
 & 0.85 & 0.897 & 0.841 & 0.950 & 0.58 \\
 & 0.80 & 0.882 & 0.823 & 0.946 & 0.56 \\
\bottomrule
\end{tabular}
\end{table}

% ------------------------------------------------------------
\subsection{Coverage per Bearing}
\label{sec:res_per_bearing}

Averages over windows can hide bearings on which the intervals fail.
Table~\ref{tab:per_bearing} reports the coverage of the full framework
for each bearing, each tested once over the five folds.

\begin{table}[htbp]
\centering
\caption{Per-bearing coverage of the full framework (+GAN+MT, CV+,
predicted stage, 90\% target), mean over five seeds. $n$: test windows
(near-failure windows).}
\label{tab:per_bearing}
\footnotesize
\renewcommand{\arraystretch}{1.15}
\begin{tabular}{lcccc}
\toprule
\textbf{Bearing} & \textbf{$n$} & \textbf{PICP all} &
\textbf{PICP near-failure} & \textbf{MPIW all} \\
\midrule
1-1 & 92 (4)      & 0.95 & 1.00 & 0.80 \\
1-2 & 130 (43)    & 0.97 & 1.00 & 0.91 \\
1-3 & 127 (45)    & 0.93 & 1.00 & 0.71 \\
1-4 & 91 (18)     & \textbf{0.76} & \textbf{0.23} & 0.26 \\
1-5 & 21 (4)      & 1.00 & 1.00 & 0.92 \\
2-1 & 460 (19)    & 0.99 & 0.96 & 0.46 \\
2-2 & 130 (17)    & 0.96 & 1.00 & 0.76 \\
2-3 & 502 (197)   & 0.92 & 1.00 & 0.75 \\
2-4 & 11 (5)      & \textbf{0.78} & \textbf{0.64} & 0.78 \\
2-5 & 308 (110)   & 0.89 & 0.99 & 0.71 \\
3-1 & 2507 (76)   & 0.99 & 1.00 & 0.45 \\
3-2 & 2465 (632)  & 0.84 & 1.00 & 0.94 \\
3-3 & 340 (10)    & 0.99 & 1.00 & 0.44 \\
3-4 & 1484 (50)   & 1.00 & 1.00 & 0.26 \\
3-5 & 83 (3)      & \textbf{0.50} & 1.00 & 0.77 \\
\bottomrule
\end{tabular}
\end{table}

Ten of the fifteen bearings are covered at 0.90 or more. Near failure,
thirteen bearings are covered at 0.96 or more. The failures are
concentrated in a few bearings. Bearing~1-4 is covered at only 0.23
near failure, with intervals of 0.26, the narrowest of the dataset
together with Bearing~3-4: the predictor is confident and wrong on
this bearing. Bearing~2-4, the shortest (11~windows), is covered at
0.78 overall and 0.64 near failure. Bearing~3-5 is covered at 0.50
overall while its three near-failure windows are covered; its errors
lie in the earlier stages. Bearing~3-2, which holds about half of all
near-failure windows, is fully covered near failure but at 0.84 overall.
Three of the four weakest bearings (1-4, 2-4 and 3-5) are among the
shortest-lived of the dataset.\\

Because the CV+ quantiles weight every calibration window equally, a
long bearing such as 3-2 contributes most of the calibration windows of
its group. A bearing-balanced variant, in which each calibration window
is weighted by the inverse number of windows of its bearing in its
group (weighted conformal prediction~\cite{tibshirani2019conformal}),
was therefore evaluated as a robustness check. For the full framework
it gives an overall coverage of $0.924 \pm 0.012$, a post-FPT coverage
of 0.893 and a near-failure coverage of 0.958, with intervals of 0.67
over the whole life. Coverage is more stable across seeds and closer to
target after the FPT, at the cost of intervals about 12\% wider. The
weakest bearings remain under-covered, so their errors are specific to
those bearings rather than a consequence of the calibration weighting.

% ------------------------------------------------------------
\subsection{Comparison with Published Results}
\label{sec:res_benchmark}

Table~\ref{tab:benchmark} places the framework among published results
on XJTU-SY. Absolute normalised RMSE values span an order of magnitude
for the same dataset, from 0.044 to 0.171, and most studies do not
state the RUL target, the scoring window or the aggregation. These
choices change the RMSE substantially: in this study the same
predictor gives 0.206 near failure, 0.275 after the FPT and 0.291 over
the whole life. A direct ranking by RMSE is therefore not meaningful.
The uncertainty-aware studies are more directly comparable. The two
bearing studies using conformal prediction report an aggregate coverage
of 0.945~\cite{wang_uncertainty-aware_2026} and 0.919 with a width of
0.654~\cite{11373138} at a 90\% target. The framework obtains 0.917
with a width of 0.60 over the whole life, and is the only one to report
coverage by stage and by bearing.

\begin{table*}[!htbp]
\centering
\caption{Reported RUL prediction results on XJTU-SY. ``NR'': not
reported in the source, which precludes strict comparability.}
\label{tab:benchmark}
\scriptsize
\renewcommand{\arraystretch}{1.15}
\setlength{\tabcolsep}{3pt}
\begin{tabular}{llcccccl}
\toprule
\textbf{Study} & \textbf{Method} & \textbf{RMSE} & \textbf{Unit} &
\textbf{Label} & \textbf{Window} & \textbf{Aggreg.} &
\textbf{Uncertainty (PICP/MPIW)} \\
\midrule
\multicolumn{8}{l}{\textit{Point-prediction methods (absolute normalised RMSE)}}\\
Jin et al.~\cite{app15179529}        & EMD--BiLSTM--Transf. & 0.044 & norm. & NR & NR & NR & none \\
Zhang et al.~\cite{11682466}         & WI-DSNAS-SS          & 0.054 & norm. & NR & NR & NR & none \\
Kumar et al.~\cite{machines14060672} & PPO                  & 0.069\textsuperscript{a} & norm. & piecew. & post-FPT & NR & none \\
An et al.~\cite{10935355}            & ILDENet              & 0.082 & norm. & NR & NR & NR & none \\
Feng et al.~\cite{FENG2026105865}    & RRDT--DCDTN          & 0.098 & norm. & NR & NR & per-bear. & none \\
Zheng et al.~\cite{10423873}         & ResNet-AE--TD3       & 0.142 & norm. & NR & NR & NR & none \\
Qiao et al.~\cite{QIAO2026112248}    & DM-SVW               & 0.16  & norm. & NR & post-FPT & per-bear. & none \\
Luo et al.~\cite{11637655}           & CNN--TCN             & 0.171 & norm. & NR & NR & per-bear. & none \\
\midrule
\multicolumn{8}{l}{\textit{Relative-improvement methods (no comparable absolute RMSE)}}\\
Lu et al.~\cite{9486650}             & GAN-LSTM             & \multicolumn{4}{l}{$-29\%$ vs LSTM (rel.)}       & time-wt. & none \\
Lu et al.~\cite{lu_joint_2022}       & HP-JT                & \multicolumn{4}{l}{$-29.4\%$ vs HP (minutes)}    & time-wt. & Gaussian layer \\
Tang et al.~\cite{10555355}          & CD-CI                & \multicolumn{4}{l}{$-52\%$ vs CNN-LSTM (rel.)}   & NR & 95\% bands (NR) \\
\midrule
\multicolumn{8}{l}{\textit{Uncertainty-aware methods}}\\
Zhang et al.~\cite{zhang_explainable_2025}     & E-GTQN      & NR & NR & NR & NR & NR & quantile+KDE: 0.940 cov. \\
Wang et al.~\cite{wang_uncertainty-aware_2026} & VAE-Transf.\ conf. & NR & norm. & NR & NR & NR & conformal: 0.945 / 0.924 \\
Bott et al.~\cite{11373138}          & Phys.-sim.+CNF conf. & 0.106\textsuperscript{b} & norm. & NR & whole & NR & conformal: 0.919 / 0.654 \\
\midrule
\rowcolor[RGB]{198,224,180}
\textbf{Proposed} & \textbf{GAN + MT BiLSTM + CV+} &
\textbf{0.206}\textsuperscript{c} & \textbf{norm.} & \textbf{piecew.} &
\textbf{stage 3} & \textbf{per-bear.} & \textbf{conformal: 0.917 / 0.60}\textsuperscript{d} \\
\bottomrule
\end{tabular}

\vspace{3pt}
\raggedright\scriptsize
\textsuperscript{a}within-bearing late-test tail only.
\textsuperscript{b}derived from the reported MSE; synthetic-trained tabular predictor.
\textsuperscript{c}near-failure stage; 0.275 post-FPT and 0.291 over the whole life.
\textsuperscript{d}whole life, 90\% target, intervals clipped to $[0,1]$; coverage by stage and by bearing in Tables~\ref{tab:intervals} and~\ref{tab:per_bearing}.
norm.: normalised. rel.: relative. conf.: conformal. cov.: coverage.
\end{table*}
